\scanpage{26}\mspage{25}
On reprend tout !

\uline{Définition-scholie} Une \uline{contrée} \inserted{$C$} est definie par un ensemble $\mathcal{L}$ (l'ens. des \uline{lieux} de la contrée), et la donnée de certains ensembles de parties de $\mathcal{L}$, appelés les \uline{figures} \inserted{(de la contrée)}.

L'ens. $\mathcal{F}$ des figures de $C$ est donc une \deleted{ensemble} partie de $\mathcal{P}(\mathcal{P}(\mathcal{L}))$ (donc $F\in \mathcal{P}(\mathcal{P}(\mathcal{P}(\mathcal{L})))$). Si \deleted{$\unsure{\mathcal{F}|F}$}$F$ est une figure, ses éléments\textsuperscript{$\times$/} (qui sont des parties de $\mathcal{L}$) s'appellent les \uline{strates} (\uline{fermées}) de la figure $F$. \unsure{Là|Lá|La} réunion des strates \deleted{ferm} de $F$ est appelée le \uline{support} de $F$, \inserted{noté $|F|$.} On pose des \uline{axiomes des contrées}, qui vont être dégagés par la suite.

\uline{Cont}~1 \quad a) La partie vide de $\mathcal{P}(\mathcal{L})$ est une \deleted{strate} figure, appelée ``\uline{figure vide}''. C'est donc la figure unique qui n'a pas de strates.

% slanted
\margin{Contenu dans \uline{Cont} 5 plus bas}
\margin{$\big\Vert$}
b) $\forall\,x\in \mathcal{L}$, \deleted{le}\inserted{il y a une} figure ayant \deleted{la} strate unique $\{x\}$. On l'appelle \uline{figure ponctuelle} definie par $x$
